% -*- coding: utf-8 -*-
% !TEX program = xelatex
\documentclass[plain,noanswer,medmath]{../../template/jnuexam}
\usepackage{../../template/kysx}

\begin{document}


\examtitle{1990}{2}


\loadproblems{1}{1990P1}%载入试卷一


\makepart{填空题}{本题满分15分，每小题3分}


% 使用试卷一，第一大题，第1小题
\useproblem{1}{1}{1}

\useproblem{1}{1}{2}

\useproblem{1}{1}{3}

\useproblem{1}{1}{4}

\useproblem{1}{1}{5}


\makepart{选择题}{本题满分15分，每小题3分}


\useproblem{1}{2}{1}

\useproblem{1}{2}{2}

\useproblem{1}{2}{3}

\useproblem{1}{2}{4}

\useproblem{1}{2}{5}


\makepart{计算题}{本题满分15分，每小题5分}


\useproblem{1}{3}{1}

\useproblem{1}{3}{2}

\useproblem{1}{3}{3}


\makepart{计算题}{本题满分18分，每小题6分}


\useproblem{1}{4}{0}%【同试卷一第四题】


\begin{problem}
求微分方程$x\ln x\dy+(y-\ln x)\dx=0$满足条件$\left.y\right|_{x=\e}=1$的特解．
\end{problem}


\begin{problem}
过点$P(1,0)$作抛物线$y=\sqrt{x-2}$的切线，该切线与上述抛物线及$x$轴围成一平面图形，
求此图形绕$x$轴旋转一周所成旋转体的体积．
\end{problem}

\makepart{}{本题满分8分}

\useproblem{1}{5}{0}%【同试卷一第五题】


\makepart{}{本题满分7分}

\useproblem{1}{6}{0}%【同试卷一第六题】


\makepart{}{本题满分6分}

\useproblem{1}{7}{0}%【同试卷一第七题】


\makepart{}{本题满分8分}

\useproblem{1}{8}{0}%【同试卷一第八题】


\makepart{}{本题满分8分}

\useproblem{1}{9}{0}%【同试卷一第九题】


\end{document}
